% !TeX root = ../cosmology_summary.tex

%% --------------------------------------------------------
\chapter{Ефект Закса-Вольфа та формування кутового спектра CMB}
\label{sec:sachs_wolfe_main}
%% --------------------------------------------------------

Кутова анізотропія температури реліктового випромінювання (CMB), яку ми спостерігаємо сьогодні (рис.~\ref{fig:planck_map}), є безпосереднім реліктовим \enquote{знімком} стану раннього Всесвіту в епоху рекомбінації ($z_{\mathrm{rec}} \approx 1100$). Флуктуації температури на небесній сфері формуються внаслідок накладання кількох фізичних процесів, що відбувалися на поверхні останнього розсіяння (LSS).


\begin{figure}[htbp]
\centering
\includegraphics[width=\textwidth]{planck_cmb_map.png}
\caption{Карта температурних анізотропій CMB від місії
         Planck~2018 \cite{planck_2018}. Кольори відображають
         відхилення від середнього $T_0 = 2.725\,\mathrm{K}$
         у діапазоні $\pm 300\,\mu\mathrm{K}$~--- спостережуваний
         відбиток первинних флуктуацій інфлатона після
         ${\sim}380\,000$ років еволюції.
         Зображення: ESA and the Planck Collaboration
         \cite{esa_planck_cmb_image}.}
\label{fig:planck_map}
\end{figure}

З погляду релятивістської кінетичної теорії, повна варіація температури у заданому напрямку $\hat{\mathbf{n}}$ описується калібровочно-інваріантним рівнянням, що враховує як внутрішні властивості плазми, так і геометрію простору-часу \cite{hu_sugiyama_1995}:
\begin{equation}\label{eq:cmb_sources_combined}
  \frac{\delta T}{\text{T}}(\hat{\mathbf{n}}) = \left( \Theta_0 + \Phi \right)_{\mathrm{rec}} + \left( \hat{\mathbf{n}} \cdot \mathbf{v}_b \right)_{\mathrm{rec}} + 2 \int_{\tau_{\mathrm{rec}}}^{\tau_0} \frac{\partial \Phi}{\partial \tau} \, d\tau.
\end{equation}
Тут перший доданок відображає \emph{ефект Закса-Вольфа} --- комбінацію внутрішньої густини фотонів $\Theta_0 \equiv \delta\rho_\gamma / 4\rho_\gamma$ та гравітаційного червоного зміщення фотонів $\Phi$ при їхньому виході з потенціальних ям. Другий доданок відповідає за доплерівський зсув частоти через пекулярний рух баріонної компоненти зі швидкістю $\mathbf{v}_b$. Третій доданок є \emph{інтегральним ефектом Закса-Вольфа} (ISW), який враховує еволюцію гравітаційних потенціалів $\dot{\Phi}$ вздовж траєкторії польоту фотонів до спостерігача внаслідок розпаду потенціалів у ранній радіаційно-домінантній епосі або під дією темної енергії на пізніх етапах.

Рівняння \eqref{eq:cmb_sources_combined} є фундаментальним, оскільки воно повністю описує локальну мікрофізику на реліктовому обрії. Проте безпосереднє зіставлення цього виразу з астрофізичними спостереженнями стикається з принциповою перешкодою: ми спроможні виміряти параметри плазми не в кожній тривимірній точці раннього Всесвіту, а лише фіксуємо двовимірну проекцію цих процесів на небеську сферу. Щобільше, конкретний візерунок гарячих і холодних плям у нашому куточку Космосу є випадковою реалізацією первинних квантових флуктуацій. Для перевірки теоретичних моделей нам потрібні не локальні значення $\delta T / T$ у конкретному напрямку, а об'єктивні статистичні характеристики, що є інваріантними відносно положення спостерігача. Саме ця потреба змушує перейти від конфігураційного простору до гармонічного аналізу сигналів на сфері.

%% --------------------------------------------------------
\section{Гармонічний аналіз та кутовий спектр потужності}
\label{subsec:harmonic_analysis}
%% --------------------------------------------------------

Оскільки спостережуваний візерунок температурних флуктуацій $\Delta T(\theta,\phi) = T(\theta,\phi) - T_\text{avg} $ фіксується на двовимірній сфері, природним математичним базисом для його статистичного опису є розклад за ортонормованими \emph{сферичними гармоніками} $Y_{\ell m}(\theta,\phi)$:
\begin{equation}\label{eq:spherical_decomp_integrated}
    \frac{\Delta T}{T}(\theta,\phi) = \sum_{\ell=0}^{\infty}\sum_{m=-\ell}^{\ell} a_{\ell m}\,Y_{\ell m}(\theta,\phi),
\end{equation}
де безрозмірні комплекснозначні коефіцієнти $a_{\ell m}$ відіграють роль амплітуд флуктуацій. Мультипольний момент $\ell$ визначає просторову частоту і є обернено пропорційним кутовому масштабу структури на небі: $\theta \approx 180^\circ/\ell$.

Для побудови теорії, вільної від вибору конкретної системи координат чи орієнтації інструментів спостереження, вводять безрозмірний \emph{кутовий спектр потужності} $C_\ell$. Припускаючи статистичну ізотропію та однорідність первинних збурень, двоточковий корелятор коефіцієнтів усереднюється за всіма можливими проекціями $m$, повністю затираючи просторові координати:
\begin{equation}
    \langle a_{\ell m} a_{\ell' m'}^* \rangle = C_\ell \, \delta_{\ell \ell'} \, \delta_{m m'}.
\end{equation}

Проте в реальних експериментах (наприклад, у космічній місії \emph{Planck}) вимірюють флуктуації безпосередньо в термодинамічних одиницях температури (мікрокельвінах, $\mu\text{K}$). Крім того, на практиці замість сирого спектра $C_\ell$ прийнято оперувати \emph{масштабованим (або згладженим) кутовим спектром потужності} $\mathcal{D}_\ell^{TT}$, який означується як:
\begin{equation}\label{eq:Dl_definition}
    \mathcal{D}_\ell^{TT} \equiv \frac{\ell(\ell+1)}{2\pi} T_0^2 \, C_\ell^{TT},
\end{equation}
де $T_0 \approx 2.7255\text{ К}$ --- середня температура реліктового фону сьогодні. Множник $\ell(\ell+1)$ має глибокий фізичний зміст: він повністю компенсує спадання амплітуди великомасштабних флуктуацій, перетворюючи чисто геометричний ефект Закса-Вольфа на плоске плато в області малих мультиполів, що значно спрощує візуальний аналіз спектра.

Прямий фізичний зв'язок між тривимірними модами первинних збурень у просторі Фур'є (що мають хвильове число $k$) та двовимірними мультиполями $\ell$ задається відстанню кутового діаметра $D_A$ до поверхні останнього розсіяння. З простої геометричної проекції плоскої хвилі довжиною $\lambda = 2\pi/k$ випливає наближене співвідношення:
\begin{equation}\label{eq:l_k_projection_main}
    \ell \approx k\cdot D_A.
\end{equation}

Строгий математичний перехід здійснюється через інтегрування за допомогою сферичних функцій Бесселя $j_\ell(kD_A)$:
\begin{equation}\label{eq:bessel_integral_main}
    C_\ell \propto \int_0^{\infty}|\Theta_k(\tau_\text{rec})|^2\cdot j_\ell^2(kD_A)\cdot k^2\,dk.
\end{equation}
Кінематична властивість функції $j_\ell(x)$ полягає в тому, що вона є практично нульовою при $x < \ell$ і зазнає різкого локального сплеску саме в околі точки $x \approx \ell$, що математично легалізує та обґрунтовує наближення \eqref{eq:l_k_projection_main}.

%% --------------------------------------------------------
\section{Зв'язок із гідродинамікою плазми та акустичні осциляції}
\label{subsec:acoustic_oscillations}
%% --------------------------------------------------------

Доки масштаб збурення перевищує розмір космологічного горизонту ($k \ll aH$), еволюція засліплена причинною незв'язаністю, і домінуючим є чисто геометричний ефект Закса-Вольфа. Для адіабатичних збурень на супергоризонтних масштабах виконується строге термодинамічне співвідношення $\Theta_0 = -\frac{2}{3}\Phi$, що після підстановки у загальне рівняння джерел \eqref{eq:cmb_sources_combined} дає відоме низькомультипольне плато Закса-Вольфа ($\ell \lesssim 30$):
\begin{equation}
  \left(\Theta_0 + \Phi\right)\Big|_{\text{супергоризонт}} = \frac{1}{3}\Phi.
\end{equation}
Введення величини $\mathcal{D}_\ell^{TT}$ згідно з формулою \eqref{eq:Dl_definition} консервує цю залежність, формуючи горизонтальний трек у лівій частині кутового спектра.

Проте, щойно збурення входить під космологічний горизонт ($k \gg aH$), градієнт внутрішнього радіаційного тиску фотонів починає чинити опір гравітаційному стисненню. Детальна мікрофізична еволюція фотон-баріонної плазми у конформно-ньютонівському калібруванні була розглянута нами у главі~\ref{sec:bao}. Зокрема, збурення фотонної густини підпорядковуються виведеному там рівнянню вимушених коливань \eqref{eq:acoustic_oscillator}:
\begin{equation}\label{eq:acoustic_oscillator_recall}
  \frac{d}{d\eta}\bigl[(1+R)\,\delta_\gamma'\bigr] + \frac{k^2}{3}\,\delta_\gamma = -\frac{k^2}{3}(1+R)\,\Phi,
\end{equation}
де безрозмірне відношення густини фотонів та баріонів $R(\eta) \equiv 3\bar{\rho}_b / 4\bar{\rho}_\gamma$ задає інерційну масу системи, а $\eta$ --- конформний час.

Для опису формування спектра CMB зручно перейти від густини $\delta_\gamma$ до монопольної компоненти збурення температури $\Theta_0 = \delta_\gamma / 4$. Якщо знехтувати повільною часовою еволюцією потенціалів ($\Phi \approx \text{const}$) та параметра $R$ в околі епохи рекомбінації, рівняння \eqref{eq:acoustic_oscillator_recall} набуває канонічного вигляду коливань гармонічного осцилятора із зовнішньою сталою сили:
\begin{equation}\label{eq:oscillator_simplified_fixed}
    \frac{d^2\Theta_0}{d\eta^2} + c_s^2\,k^2\,\Theta_0 = -\frac{k^2}{3}\,\Phi,
\end{equation}
де конформна швидкість звуку $c_s = 1/\sqrt{3(1+R)}$ повністю узгоджується з раніше виведеним виразом \eqref{eq:sound_speed}.

У безсиловому наближенні ($\Phi = 0$), розв'язок рівняння \eqref{eq:oscillator_simplified_fixed} для адіабатичних початкових умов задається гармонічними коливаннями: $\Theta_k(\eta) \approx \Theta_k(0)\cos(kc_s\eta)$. Миттєва дезактивація вільних електронів під час рекомбінації працює як термодинамічний \enquote{стоп-кадр}, що фіксує фазу коливань на моменті конформного часу $\eta_\text{rec}$. Оскільки внесок у кутову анізотропію на малих масштабах визначається саме цією замороженою хвилею, кутовий спектр потужності стає пропорційним її квадрату:
\begin{equation}
    C_\ell \propto \cos^2\!\left(kc_s\eta_\text{rec}\right).
\end{equation}
Екстремуми цієї функції при $kc_s\eta_\text{rec} = \pi, 2\pi, 3\pi,\ldots$ формують періодичну структуру акустичних піків у спектрі анізотропії, які чітко реєструються на графіку $\mathcal{D}_\ell^{TT}$.

Важливо наголосити, що підгоризонтні звукові хвилі не затерли первинний розподіл густини, згенерований інфлатоном, а лише \emph{модулювали} його. Оскільки темна матерія не взаємодіє з випромінюванням, її гравітаційні потенціальні ями залишалися статичними та нерухомими, виступаючи стабільною координатною сіткою. Баріон-фотонний газ коливався відносно цих застиглих центрів, де $\Theta_k(0)$ є первісним квантовим імпульсом інфлатона. Фізичною аналогією тут є глибокий рельєфний малюнок на папері, який потрапив под піщану бурю: вітер не згладжує борозни, а навпаки --- пісок засипається в заглиблення, контрастно підкреслюючи первинну геометрію середовища.

%% --------------------------------------------------------
\section{Ефект баріонного завантаження та порушення симетрії}
\label{subsec:baryon_loading_effects}
%% --------------------------------------------------------

Наявність важкої баріонної компоненти, що описується параметром $R$, відображається на внутрішній симетрії фотонних осциляцій. Якщо повернутися до точного розв'язку неоднорідного рівняння \eqref{eq:oscillator_simplified_fixed} зі збереженням гравітаційного члена у правій частині, виявляється, що присутність баріонів зміщує положення динамічної рівноваги (нуля) коливань осцилятора. Математично, загальний розв'язок для монополя температури набуває вигляду:
\begin{equation}\label{eq:shifted_oscillator_final_fixed}
    \Theta_0(\eta_\text{rec}) = \Theta_0(0)\cdot\cos\!\left(kc_s\eta_\text{rec}\right) - (1+R)\Phi.
\end{equation}

З урахуванням ефекту Закса-Вольфа, повне ефективне збурення, яке виривається з потенціальної ями, визначається комбінацією $\Theta_0 + \Phi$. Підставляючи сюди розв'язок \eqref{eq:shifted_oscillator_final_fixed}, знаходимо:
\begin{equation}\label{eq:baryon_loading_shift_result}
    \left( \Theta_0 + \Phi \right)_{\mathrm{rec}} = \Theta_0(0)\cdot\cos\!\left(kc_s\eta_\text{rec}\right) - R\Phi.
\end{equation}

Безрозмірний зсув нуля $\Delta \equiv -R\Phi / \Theta_0(0)$ створює суттєву асиметрію між парними та непарними акустичними піками у спектрі потужності $\mathcal{D}_\ell^{TT}$:
\begin{itemize}
    \item \textbf{Непарні піки (перший, третій тощо)} відповідають фазі максимального гравітаційного стиснення плазми всередині потенціальних ям темної матерії. Гравітаційне поле баріонів діє синергетично з колапсом, підсилюючи амплідуду ефективного збурення до величини $(1+\Delta)$. Відповідно, кутова потужність пропорційна квадрату цього значення:
    \begin{equation}
        \text{Потужність непарних піків} \propto (1+\Delta)^2.
    \end{equation}

    \item \textbf{Парні піки (другий, четвертий тощо)} фіксують протилежну фазу --- момент максимального пружного розширення (радіаційного розрідження). У цій точці фотонний тиск намагається викинути речовину з ями, але гравітаційний потенціал баріонів утримує її, послаблюючи чисту амплітуду до $(-1+\Delta)$. Для кутового спектра потужності це дає:
    \begin{equation}
        \text{Потужність парних піків} \propto (-1+\Delta)^2 = (1-\Delta)^2.
    \end{equation}
\end{itemize}

Оскільки і метричний потенціал $\Phi$, і початкове збурення плазми $\Theta_0(0)$ народилися з єдиного первинного скалярного поля, їхнє відношення є безрозмірною константою порядку одиниці. Якби у Всесвіті були відсутні баріони ($R=0 \implies \Delta = 0$), осцилятор коливався б симетрично відносно нуля, і перший та другий піки мали б абсолютно однакову висоту. Отже, прецизійне вимірювання відносної висоти перших двох піків спектра $\mathcal{D}_\ell^{TT}$ дозволяє астрофізикам з високою точністю \enquote{зважити} баріонний вміст Всесвіту, відокремивши його від внеску холодної темної матерії (див. рис.~\ref{fig:cmb_tt}).

\begin{figure}[htbp]
\centering
\begin{tikzpicture}
\begin{semilogxaxis} [
    width   = \linewidth,
    height  = \linewidth,
    xlabel  = {Мультипольний момент $\ell$},
    ylabel  = {$\mathcal{D}_\ell^{TT}\ [\mu\mathrm{K}^2]$},
    yticklabel style={/pgf/number format/1000 sep={\,}},
    xmin    = 2,   xmax = 2600,
    ymin    = 0,   ymax = 6800,
    % --- РЯТУЄМО ОСЬ Х ПІД \linewidth ---
    xtick = {2, 10, 30, 100, 500, 1000, 1500, 2000, 2500},
    % Використовуємо компактний науковий формат, щоб мітки не зливалися на вузькому графіку:
    xticklabels = {2, 10, 30, 100, 500, , 1500},
    % Трохи зменшуємо шрифт міток для осі Х та зміщуємо вниз, щоб сітка не заважала
    ticklabel style = {font=\footnotesize, yshift=-1pt},
    % Твоя густа логарифмічна сітка
    minor xtick = {3,4,5,6,7,8,9,20,40,50,60,70,80,90,200,300,400,600,700,800,900,1100,1200,1300,1400,1600,1700,1800,1900,2100,2200,2300,2400},
    ytick   = {0, 1000, 2000, 3000, 4000, 5000, 6000},
    grid   = both,
    grid style = {line width=0.3pt, draw=gray!25},
    major grid style = {line width=0.5pt, draw=gray!40},
    label style      = {font=\small},
    legend style = {
        at={(0.6,0.97)}, anchor=north east,
        font=\small, fill=white, fill opacity=0.9,
        draw=gray!60, inner sep=4pt, row sep=1pt
    },
    clip = true,
]

    %--- LambdaCDM Best-fit Curve ---
    \addplot[
        color     = theoryred,
        width     = 1.6pt,
        smooth,
    ] table[x=ell, y=Dl] {\theorydata};
    \addlegendentry{$\Lambda$CDM (CAMB, Planck\,2018 best-fit)}

    %--- Planck 2018 data with error bars ---
    \addplot[
        only marks,
        mark       = *,
        mark size  = 1.4pt,
        color      = planckblue,
        error bars/.cd,
            y dir  = both,
            y explicit,
    ] table[
        x    = ell,
        y    = Dl,
        y error minus = errm,
        y error plus  = errp
    ] {\planckdata};
    \addlegendentry{Planck\,2018 (збіновані дані TT)}

    %--- Annotate acoustic peaks ---
    \node[font=\tiny, text=gray!70, above] at (axis cs:220,5732) {1-й пік};
    \node[font=\tiny, text=gray!70, above] at (axis cs:540,2592) {2-й пік};
    \node[font=\tiny, text=gray!70, above] at (axis cs:810,2540) {3-й пік};

    %--- Sachs-Wolfe plateau annotation ---
    \draw[->, gray!60, thin] (axis cs:6,1300) -- (axis cs:30,1100);
    \node[font=\tiny, text=gray!70, anchor=south, align=center] at (axis cs:12,1350) {Плато\\Закса-Вольфа};

  \end{semilogxaxis}
\end{tikzpicture}
\caption{Кутовий спектр потужності температурної анізотропії реліктового випромінювання за даними космічної місії \emph{Planck 2018}. Суцільна червона лінія відповідає найкращому фіту стандартної космологічної моделі $\Lambda$CDM.}
\label{fig:cmb_tt}
\end{figure}